\documentclass[11pt]{article}
\usepackage[T1]{fontenc}
\usepackage[utf8]{inputenc}
\usepackage{amsmath,amssymb}
\usepackage[margin=1in]{geometry}
\usepackage{graphicx}
\usepackage{booktabs}
\usepackage{float}
\usepackage[colorlinks=true,linkcolor=blue,citecolor=blue]{hyperref}

\title{Conditional Mutual Information and Petz Recovery\\
in a $\mathbb{Z}_2$ Lattice Gauge Ground State}
\author{Lluis Eriksson\\
\small Independent Researcher\\
\small \texttt{lluiseriksson@gmail.com}}
\date{January 2026 (v1); July 2026 (v2)}

\begin{document}
\maketitle

\begin{abstract}
We study approximate quantum Markov structure in a $\mathbb{Z}_2$
lattice gauge ground state using the conditional mutual information
(CMI) $I(A : C \mid B(w))$ and the performance of Petz recovery across a
family of tripartitions $(A, B(w), C)$ parameterized by a buffer width
$w$. We consider a $2\times4$ plaquette lattice with open boundaries and
qubits on links, restricted to a gauge-invariant (Gauss-law) physical
sector, at coupling $g = 1.0$. For each $w$ we compute reduced density
matrices, the entropies entering the CMI, and a Petz-recovered state
$\sigma_{ABC} = (\mathrm{id}_A \otimes
\mathcal{R}^{\mathrm{Petz}}_{B\to BC})(\rho_{AB})$, reporting fidelity
$F(\rho_{ABC}, \sigma_{ABC})$ via the recovery error $E_{\mathrm{rec}}(w)
= -\log F$. The Overleaf project includes the plot, a formatted table,
raw CSV outputs, and a hash-based manifest; the appendix typesets raw
artifacts. We also report numerical cross-checks (dense vs.\ low-rank
method agreement and trace stability) to support validity. \emph{v2 (no
v1 number is changed):} the star-operator definition is corrected ---
with the Hamiltonian convention used here (single-link $Z$ terms), the
Gauss stars must be $G_s = \prod_{\ell \in \mathrm{star}(s)} Z_\ell$;
v1's printed $\prod X_\ell$ anticommutes with the $Z_\ell$ terms of $H$
(the code used the consistent convention: an independent reconstruction
from the manifest alone reproduces the CSV ground energy to $7\times
10^{-15}$ and every CMI of Table~\ref{tab:data} to machine precision);
two interpretive remarks are added --- the CMI rise at $w = 2$ tracks the
shrinking traced-out complement ($|D|: 8 \to 2 \to 0$), so the profile
is not a shielding-decay curve, and the $w = 2 \approx w = 3$ plateau is
the buffer-saturation identity of the companion 2601.0050 (v2); the
apparent Petz-over-CMI excess at $w = 1$ ($E_{\mathrm{rec}} > I$) is
shown to be the $\delta = 10^{-6}$ regularization floor --- regenerating
with $\delta = 10^{-12}$ restores $E_{\mathrm{rec}} \le I$ at every $w$
in the regenerated dataset; and a verification suite replicates the full
pipeline from the manifest data alone.
\end{abstract}

\section{Introduction}

The conditional mutual information (CMI) quantifies residual
correlations between $A$ and $C$ conditioned on $B$:
\[
I(A : C \mid B) = S(\rho_{AB}) + S(\rho_{BC}) - S(\rho_B) -
S(\rho_{ABC}),
\]
where $S(\rho) = -\mathrm{tr}(\rho\log\rho)$. Small CMI is
characteristic of (approximate) quantum Markov chains and is closely
tied to the possibility of reconstructing $\rho_{ABC}$ from $\rho_{AB}$
by a recovery channel acting only on $B$. A canonical candidate is the
Petz map; more generally, several recovery constructions yield
quantitative links between CMI and achievable reconstruction fidelity
\cite{Petz, FR, SBT}.

Lattice gauge models provide a concrete setting to test these ideas
under gauge constraints. In this short reproducible note we compute CMI
and Petz recovery across a controlled buffer geometry in a
$\mathbb{Z}_2$ gauge model, and provide all numerical artifacts for
verification.

\section{Model and gauge-invariant sector}

We place qubits on links $\ell$ of a square lattice with open boundary
conditions. Let $X_\ell$ and $Z_\ell$ denote Pauli operators on link
$\ell$. We use the Hamiltonian convention implemented in the
accompanying code:
\[
H(g) = -g \sum_\ell Z_\ell - \frac{1}{g} \sum_p X_p, \qquad
X_p = \prod_{\ell \in \partial p} X_\ell,
\]
where $p$ runs over plaquettes and $\partial p$ denotes the boundary
links of plaquette $p$. We work at $g = 1.0$ on a $2\times4$ plaquette
lattice.

Gauge invariance is enforced by restricting to the physical (Gauss-law)
subspace generated by star operators
\[
G_s = \prod_{\ell \in \mathrm{star}(s)} Z_\ell, \qquad G_s\psi = \psi
\ \text{for all vertices } s.
\]
In this convention one has $[H(g), G_s] = 0$ for all vertices $s$, so
the dynamics preserves the $G_s = +1$ sector. We compute the ground
state $\psi_0$ in this physical sector and set $\rho =
\psi_0\psi_0^\dagger$.

\textbf{v2 correction (star operators).} v1 printed the stars as
$\prod_{\ell\in\mathrm{star}(s)} X_\ell$. With the single-link $Z_\ell$
terms of $H(g)$ above this is inconsistent: each $Z_\ell$ shares exactly
one link with the two stars at its endpoints and therefore
\emph{anticommutes} with them, so $[H, G_s] \ne 0$ as printed. The
consistent choice for this Hamiltonian convention --- and the one
actually realized in the numerical pipeline --- is $G_s = \prod Z_\ell$
as displayed above; the $G_s = +1$ sector is then the cycle space
generated by plaquette flips, of dimension $2^{N_p} = 256$, matching
\texttt{dim\_phys} in the manifest. (Equivalently, one may swap $X
\leftrightarrow Z$ everywhere, recovering the standard Kogut convention;
the two readings are related by a global Hadamard.) The v2 verification
suite confirms that this reconstruction reproduces the CSV ground energy
$E_0 = -22.997180401986128$ to $7\times10^{-15}$ and every entry of
Table~\ref{tab:data} (see Note added).

\section{Regions $A$, $B(w)$, $C$ and buffer construction}

Subsystems are defined as subsets of links; reduced states are obtained
by tracing out complement links. We fix disjoint link sets $A$ and $C$
(``terminals'') and define a separating wall $W$ between them, chosen as
a brute-force minimum vertex cut in the vertex-adjacency graph used by
the pipeline.

The pipeline uses two link-adjacency graphs: \emph{vertex-adjacency}
(links adjacent if they share a vertex), used to construct and verify a
separating wall $W$; and \emph{plaquette-adjacency} (links adjacent if
they belong to a common plaquette), used to grow the buffer.

Let $d_\square(\ell, W)$ be the graph distance from link $\ell$ to the
wall $W$ in the plaquette-adjacency graph. For integer $w \ge 0$ we
define the buffer region
\[
B(w) = \{\ell : d_\square(\ell, W) \le w\} \setminus (A \cup C),
\]
ensuring $A$, $B(w)$, and $C$ are disjoint by construction. The exact
link indexing and the sets for each $w$ are recorded in
\texttt{repro\_manifest.json} and \texttt{z2\_petz\_results.csv}.

\section{Methods: CMI and Petz recovery}

For each $w$ we form reduced density matrices $\rho_{AB}$, $\rho_{BC}$,
$\rho_B$, and $\rho_{ABC}$ by partial traces, and compute entropies from
eigenvalues using a small eigenvalue cutoff to avoid numerical $\log(0)$
instabilities (the cutoff value is recorded in the manifest). The
implementation is hybrid: we use dense linear algebra when $|ABC|$ is
small, and a low-rank (LR+OB) pipeline for larger buffers; both methods
are cross-checked in the overlap regime. We evaluate Petz recovery only
for instances where the constructed wall separates $A$ and $C$ (recorded
as \texttt{SEP} in the CSV).

To quantify local reconstructibility we apply the Petz recovery map:
\[
\mathcal{R}^{\mathrm{Petz}}_{B\to BC}(X) = \rho_{BC}^{1/2}\,
\rho_B^{-1/2}\, X\, \rho_B^{-1/2}\, \rho_{BC}^{1/2}.
\]
Numerically, we stabilize the inverse square root by a fixed
regularization $\rho_B^{-1/2} \approx (\rho_B + \delta I)^{-1/2}$,
$\delta = 10^{-6}$, as recorded in the manifest (parameter
\texttt{DELTA\_PETZ}). The recovered state is $\sigma_{ABC} =
(\mathrm{id}_A \otimes \mathcal{R}^{\mathrm{Petz}}_{B\to
BC})(\rho_{AB})$, and we explicitly renormalize $\sigma_{ABC} \leftarrow
\sigma_{ABC}/\mathrm{tr}(\sigma_{ABC})$ before computing fidelity. We
additionally symmetrize $\sigma_{ABC} \leftarrow (\sigma_{ABC} +
\sigma_{ABC}^\dagger)/2$ to suppress tiny non-Hermitian numerical noise.
We verify $\sigma_{ABC} \succeq 0$ up to numerical tolerance before
evaluating the Uhlmann fidelity
\[
F(\rho_{ABC}, \sigma_{ABC}) = \Bigl(\mathrm{tr}\sqrt{
\sqrt{\rho_{ABC}}\, \sigma_{ABC} \sqrt{\rho_{ABC}}}\Bigr)^2
\]
and the recovery error $E_{\mathrm{rec}}(w) = -\log F$ (natural
logarithm, nats).

\section{Results}

Figure~\ref{fig:money} shows $I(A : C \mid B(w))$ together with
$E_{\mathrm{rec}}(w)$. The CSV records whether each point was computed
using the dense method or the LR+OB method (field
\texttt{petz\_method}). Table~\ref{tab:data} reports the numerical
values, including $1 - F$. The main qualitative observation in this
instance is that Petz recovery remains high-fidelity (with recovery
error of order $10^{-5}$) over the tested buffer widths $w \in \{0, 1,
2, 3\}$, consistent with an approximately Markovian structure for this
geometry at $g = 1.0$.

\begin{figure}[h!]
\centering
\includegraphics[width=0.8\linewidth]{fig1.png}
\caption{$\mathbb{Z}_2$ model on a $2\times4$ plaquette lattice at $g =
1.0$. Conditional mutual information $I(A : C \mid B(w))$ and Petz
recovery error $E_{\mathrm{rec}}(w) = -\log F$ versus buffer width $w$.}
\label{fig:money}
\end{figure}

\begin{table}[H]
\centering
\caption{$\mathbb{Z}_2$ model on a $2\times4$ plaquette lattice at $g =
1.0$. Region sizes and values of $I(A : C \mid B(w))$ and $1 - F$.}
\label{tab:data}
\begin{tabular}{c c c c c}
\toprule
$w$ & $|A|$ & $|B|$ & $I(A : C \mid B)$ & $1 - F$ \\
\midrule
0 & 4 & 3 & 1.627985e-07 & 8.898586e-09 \\
1 & 4 & 6 & 1.659899e-07 & 2.422318e-07 \\
2 & 4 & 12 & 5.782748e-05 & 2.589104e-05 \\
3 & 4 & 14 & 5.782759e-05 & 2.588882e-05 \\
\bottomrule
\end{tabular}
\end{table}

\textbf{Interpretive remarks (v2).} Two features of Table~\ref{tab:data}
deserve explicit framing. \emph{(i) The profile is not a shielding-decay
curve.} With $A$ and $C$ fixed and $B(w)$ growing, the traced-out
complement $D$ shrinks ($|D| = 11, 8, 2, 0$ for $w = 0..3$); the
$\times 350$ CMI rise at $w = 2$ coincides exactly with $|D|$ dropping
from $8$ to $2$ links. This is the complement-shrinkage confound
documented for finite-size pipelines in the companions 2601.0038 and
2601.0040 (v2): $w$ simultaneously changes the buffer and the discarded
environment, so monotone decay in $w$ should not be expected, and small
CMI at small $w$ partly reflects a large traced-out region.
\emph{(ii) The $w = 2 \approx w = 3$ plateau is saturation.} At $w = 3$
the buffer is the full complement of $A \cup C$ and the global state is
pure; the near-equality $I(w{=}2) \simeq I(w{=}3)$ (relative difference
$2\times10^{-6}$) is the buffer-saturation identity proved in the
companion 2601.0050 (v2). \emph{(iii) The $w = 1$ Petz-over-CMI excess is a
$\delta$ floor.} In Table~\ref{tab:data}, $E_{\mathrm{rec}}(1) =
2.42\times10^{-7} > I(1) = 1.66\times10^{-7}$, numerically above the CMI scale (Fawzi--Renner guarantees
$E_{\mathrm{rec}} \le I$ for \emph{some} recovery channel, not for the
explicit regularized Petz map pointwise, so this is a diagnostic, not a
violation). Regenerating the
dataset with $\delta = 10^{-12}$ gives $E_{\mathrm{rec}}(1) =
2.3\times10^{-8} \le I(1)$, and $E_{\mathrm{rec}} \le I$ at every $w$ in
the regenerated dataset: the excess is the $\delta = 10^{-6}$
regularization floor, not a property of the Petz map here.

\section{Validation checks}

We include lightweight numerical checks to support reproducibility:
\begin{itemize}
\item \textbf{Method cross-check.} At the smallest buffer width where
both implementations are feasible, dense evaluation versus the LR+OB
pipeline agrees at the level $|\Delta E_{\mathrm{rec}}| \approx 5.394
\times 10^{-8}$, and the trace agreement satisfies $|\Delta\mathrm{tr}|
\approx 5.6\times10^{-16}$ (see CSV for the corresponding rows).
\item \textbf{Trace stability.} The recovered state remains close to
normalized; for example we observe $\mathrm{tr}(\sigma) \approx
0.9999855$ in the reported run (see CSV).
\end{itemize}

\section{Scope and outlook}

This paper reports one system size ($2\times4$) and one coupling ($g =
1.0$). Immediate extensions include scanning $g$ across regimes,
finite-size scaling, and comparing Petz recovery to alternative recovery
maps and precision controls. Within the series, the coupling sweeps and
confinement diagnostics of 2601.0044, 2601.0047 and 2601.0051 (all v2)
are the natural continuation; the fixed-terminal geometry here
complements their collar-decay geometries.

\section{Reproducibility artifacts}

This Overleaf project expects the following files in the project root:
\texttt{fig\_z2\_petz\_recovery.png}, \texttt{tab\_z2\_petz\_data.tex},
\texttt{z2\_petz\_results.csv}, \texttt{repro\_manifest.json}. The
manifest records run metadata (model and numerical parameters, cutoffs)
together with file hashes to enable integrity checks of the artifacts.
\emph{v2:} the verification suite
(\texttt{verification/2601-0111/verify\_2601\_0111.py} in the series
repository) reconstructs the entire pipeline from the manifest fields
alone (\texttt{A\_patch}, \texttt{C\_patch}, \texttt{wall},
\texttt{W\_LIST}, \texttt{DELTA\_PETZ}) and replicates
Table~\ref{tab:data}; see the Note added.

\section*{Note added in v2}
No v1 number is changed; the CSV, manifest, table and figure stand. v2:
(1) corrects the star-operator definition ($G_s = \prod Z_\ell$; v1's
printed $\prod X_\ell$ anticommutes with the $Z_\ell$ terms of $H$, so
$[H, G_s] \ne 0$ as printed --- a typo, since the pipeline used the
consistent convention); (2) adds the interpretive remarks above:
complement-shrinkage confound (the CMI rise at $w = 2$ tracks $|D|: 8
\to 2$), buffer saturation at $w = 3$ (2601.0050 v2), and the
$\delta$-floor diagnosis of the $w = 1$ Petz-over-CMI excess ($\delta = 10^{-12}$
restores $E_{\mathrm{rec}} \le I$ at every $w$ in the regenerated
dataset); (3) adds series positioning (v1 cited no companion papers);
(4) adds a verification suite that reconstructs the lattice, the v1 link
indexing, the buffers ($|B(w)| = 3, 6, 12, 14$ reproduced exactly), the
cycle-space Gauss sector ($\dim = 256$), the ground state ($E_0$ matched
to $7\times10^{-15}$), all four CMI values to machine precision
($<10^{-13}$), and $1 - F$ to $10^{-6}$ relative on the LR+OB rows of
Table~\ref{tab:data} ($w = 1, 2, 3$); at $w = 0$ (v1's dense row) the
suite's low-rank method differs from v1's dense value by
$1.2\times10^{-8}$, within the $5.4\times10^{-8}$ dense-vs-LR+OB method
spread that v1 itself reports; the recomputed $\mathrm{tr}(\sigma)$
values match the CSV to all printed digits.

\begin{thebibliography}{9}
\bibitem{Petz} D. Petz, \emph{Sufficient subalgebras and the relative
entropy of states of a von Neumann algebra}, Commun. Math. Phys.
\textbf{105} (1986) 123--131.
\bibitem{FR} O. Fawzi and R. Renner, \emph{Quantum conditional mutual
information and approximate Markov chains}, Commun. Math. Phys.
\textbf{340} (2015) 575--611.
\bibitem{SBT} D. Sutter, M. Berta, and M. Tomamichel,
\emph{Multivariate trace inequalities}, Commun. Math. Phys. \textbf{352}
(2017) 37--58.
\bibitem{Wegner} F. J. Wegner, \emph{Duality in generalized Ising models
and phase transitions without local order parameters}, J. Math. Phys.
\textbf{12} (1971) 2259.
\bibitem{Kogut} J. B. Kogut, \emph{An introduction to lattice gauge
theory and spin systems}, Rev. Mod. Phys. \textbf{51} (1979) 659.
\bibitem{Kitaev} A. Y. Kitaev, \emph{Fault-tolerant quantum computation
by anyons}, Ann. Phys. \textbf{303} (2003) 2--30.
\bibitem{S0050} L. Eriksson, ai.viXra:2601.0050 (2026; v2 2026).
\bibitem{S0038} L. Eriksson, ai.viXra:2601.0038 (2026; v2 2026).
\bibitem{S0040} L. Eriksson, ai.viXra:2601.0040 (2026; v2 2026).
\bibitem{S0044} L. Eriksson, ai.viXra:2601.0044 (2026; v2 2026).
\bibitem{S0047} L. Eriksson, ai.viXra:2601.0047 (2026; v2 2026).
\bibitem{S0051} L. Eriksson, ai.viXra:2601.0051 (2026; v2 2026).
\end{thebibliography}

\appendix

\section{Reproducibility appendix}

\subsection{Raw data (CSV): \texttt{z2\_petz\_results.csv}}

{\scriptsize
\begin{verbatim}
Nx,Ny,g,E0_phys,w,nA,nB,nC,sep,nABC,CMI,Erec,F,1-F,tr_sigma,wall_size,petz_method
2,4,1.0,-22.997180401986128,0,4,3,4,SEP,11,1.6279854714307262e-07,8.898585574692225e
-09,0.9999999911014144,8.898585646122115e-09,0.9999960024027247,3,[DENSE]
2,4,1.0,-22.997180401986128,1,4,6,4,SEP,14,1.6598989827087962e-07,2.4223181230346393e
-07,0.9999997577682169,2.4223178307636317e-07,0.9999853324155468,3,[LR+OB rABC=256 rB=32]
2,4,1.0,-22.997180401986128,2,4,12,4,SEP,20,5.782748071088417e-05,2.5891375884829603e
-05,0.9999741089592938,2.5891040706160773e-05,0.9999855315493759,3,[LR+OB rABC=4 rB=48]
2,4,1.0,-22.997180401986128,3,4,14,4,SEP,22,5.782758913397412e-05,2.58891519139754e
-05,0.9999741111832071,2.5888816792885017e-05,0.9999855315493761,3,[LR+OB rABC=1 rB=48]
\end{verbatim}
}

\subsection{Manifest (hashes + metadata): \texttt{repro\_manifest.json}}

Key fields (full file in the repository): \texttt{A\_patch} $= [0, 2,
10, 11]$; \texttt{C\_patch} $= [6, 8, 19, 20]$; \texttt{wall} $= [13,
14, 15]$ (\texttt{wall\_size} $= 3$); \texttt{W\_LIST} $= [0, 1, 2,
3]$; \texttt{DELTA\_PETZ} $= 10^{-6}$; \texttt{DENSE\_MAX\_NABC} $=
12$; \texttt{Nlinks} $= 22$, \texttt{Np} $= 8$, \texttt{Nv} $= 15$,
\texttt{dim\_phys} $= 256$; SHA256 hashes of
\texttt{fig\_z2\_petz\_recovery.png} (\texttt{17d33a18\ldots}),
\texttt{tab\_z2\_petz\_data.tex} (\texttt{d07e016d\ldots}),
\texttt{z2\_petz\_results.csv} (\texttt{710f9a48\ldots}).

\end{document}
